% Primary and supplementary source basis: :contentReference[oaicite:7]{index=7} :contentReference[oaicite:8]{index=8}

\section{Complex Functions}

\begin{conceptbox}
    A \textbf{complex function} is a mapping
    \[
        f:\mathbb{C}\to\mathbb{C}.
    \]
    If
    \[
        z=x+iy,
    \]
    then a complex function can be written as
    \[
        w=f(z)=u(x,y)+iv(x,y),
    \]
    where $u$ and $v$ are real-valued functions of two real variables.
\end{conceptbox}

\subsection*{Motivation / Intuition}

A real function maps a line to a line. A complex function maps a plane to a plane. That means a complex function can move, rotate, stretch, fold, or invert regions of the plane.

This is why visualizing complex functions often involves studying how curves and regions in the $z$-plane map into the $w$-plane.

\subsection*{Geometric Interpretation}

Useful curves in the complex plane include:
\[
    z=x+iy_0 \quad \text{(horizontal line)},
\]
\[
    z=x_0+iy \quad \text{(vertical line)},
\]
\[
    z=re^{i\theta_0}, \quad r\ge 0 \quad \text{(ray)},
\]
\[
    |z-c|=R \quad \text{(circle)}.
\]

A complex mapping may turn these simple curves into other lines, circles, spirals, or more complicated shapes.

\subsection*{Key Properties and Observations}

\begin{conceptbox}
    If
    \[
        f(z)=a+bz,
    \]
    then:
    \begin{itemize}
        \item adding $a$ translates points,
        \item multiplying by a real $b$ dilates the plane,
        \item multiplying by a complex number $b=\rho e^{i\phi}$ both dilates by $\rho$ and rotates by $\phi$.
    \end{itemize}
\end{conceptbox}

\begin{conceptbox}
    If
    \[
        f(z)=z^2,
    \]
    then
    \[
        z=re^{i\theta}
        \quad \Longrightarrow \quad
        w=r^2e^{i2\theta}.
    \]
    So the magnitude is squared and the argument is doubled.
\end{conceptbox}

\begin{conceptbox}
    If
    \[
        f(z)=\frac{1}{z},
    \]
    then
    \[
        z=re^{i\theta}
        \quad \Longrightarrow \quad
        w=\frac{1}{r}e^{-i\theta}.
    \]
    So magnitudes are inverted and angles change sign.
\end{conceptbox}

\begin{examplebox}
    \textbf{Example 1: Translation and rotation}

    Consider
    \[
        f(z)=2e^{i\pi/6}z+1.
    \]

    Let
    \[
        z=re^{i\theta}.
    \]

    Then
    \[
        f(z)=2re^{i(\theta+\pi/6)}+1.
    \]

    Interpretation:
    \begin{enumerate}
        \item multiply distances from the origin by $2$,
        \item rotate every point by $\pi/6$,
        \item then translate the entire image one unit to the right.
    \end{enumerate}
\end{examplebox}

\begin{examplebox}
    \textbf{Example 2: Mapping under $w=z^2$}

    Let
    \[
        z=e^{i\theta}, \qquad 0\le \theta \le \pi.
    \]

    These points trace the upper semicircle of the unit circle.

    Under
    \[
        w=z^2,
    \]
    we get
    \[
        w=e^{i2\theta}.
    \]

    As $\theta$ goes from $0$ to $\pi$, the angle $2\theta$ goes from $0$ to $2\pi$.

    Therefore, the upper semicircle maps onto the entire unit circle.
\end{examplebox}

\begin{examplebox}
    \textbf{Example 3: Write $1/z$ in rectangular components}

    Let
    \[
        w=\frac{1}{z}, \qquad z=x+iy.
    \]

    Then
    \[
        w=\frac{1}{x+iy}\cdot \frac{x-iy}{x-iy}
        =
        \frac{x-iy}{x^2+y^2}.
    \]

    Hence
    \[
        u(x,y)=\frac{x}{x^2+y^2},
        \qquad
        v(x,y)= -\frac{y}{x^2+y^2}.
    \]
\end{examplebox}

\begin{noteBox}
    Complex mappings are fundamental in engineering because they connect geometry with system behavior. For example, mappings between the $s$-plane and the $z$-plane are central in control systems and digital signal processing.
\end{noteBox}

\subsection*{Common Mistakes / Notes}

\begin{itemize}
    \item A complex function is not just a pair of unrelated real functions; $u$ and $v$ are linked when differentiability is required.
    \item The map $z^2$ is not one-to-one on the whole plane.
    \item The map $1/z$ is undefined at $z=0$ and has singular behavior there.
\end{itemize}